\chapter{Major minerals and electrolytes}

\section{Learning objectives}
Describe the distribution and regulation of calcium, phosphorus, magnesium, sodium, potassium, chloride, and sulfur; relate electrolytes to neuromuscular and fluid physiology; and identify circumstances in which supplementation or salt substitution can be hazardous.

\section{Calcium and phosphorus}
The clinical information in this section is summarized from current nutrient reference material and should not be used as an individualized replacement protocol. \cite{odsList}
Calcium supports bone, teeth, muscle contraction, nerve signaling, and blood clotting. Dairy foods, calcium-set tofu, fortified beverages, canned fish with bones, and some greens are sources. Vitamin D and adequate protein support bone health. Too little calcium over time can contribute to bone loss; excess supplemental calcium may cause constipation, kidney stones, or medication interactions. Phosphorus is abundant in protein foods and many additives; deficiency is uncommon except in particular illnesses or treatments.

Approximately 99\% of body calcium is stored in bone, which functions as a reservoir. Ionized calcium is tightly regulated by parathyroid hormone, vitamin D, kidney function, and intestinal absorption. Serum calcium is therefore a poor measure of total bone stores. Osteoporosis is diagnosed by fracture risk and bone-density assessment, not by a low dietary intake alone. \cite{odsCa}

\section{Magnesium}
Magnesium participates in more than 300 enzyme systems, including energy production, protein synthesis, muscle and nerve function, glucose regulation, and blood-pressure regulation. Legumes, nuts, seeds, whole grains, and leafy greens are rich sources. Diarrhea is common with some supplemental forms; severe excess is primarily a concern with kidney impairment. \cite{odsMg}

Total serum magnesium may remain normal despite intracellular depletion. Risk increases with chronic diarrhea, malabsorption, alcohol use disorder, uncontrolled diabetes, and medications that increase renal loss. Severe deficiency can cause weakness, tremor, seizures, arrhythmias, hypokalemia, or hypocalcemia.

\section{Sodium, potassium, and chloride}
Sodium and chloride help maintain fluid balance and nerve impulses. Most sodium comes from processed and restaurant foods rather than a salt shaker. Potassium is found in beans, potatoes, fruit, vegetables, dairy, and fish. For many adults, replacing highly processed salty foods with potassium-rich foods supports blood pressure, but potassium supplements or salt substitutes can be dangerous in kidney disease or with medicines such as ACE inhibitors, ARBs, or potassium-sparing diuretics.

\section{Sulfur}
Sulfur is part of amino acids and several biologically active molecules. Healthy diets provide it through protein foods and sulfur-containing vegetables such as onions, garlic, cabbage, and broccoli. There is no established general-purpose sulfur supplement requirement.

\begin{table}[ht]
\centering
\caption{Electrolyte safety principle}
\begin{tabular}{p{0.37\textwidth}p{0.48\textwidth}}
\toprule
Situation & Why professional guidance matters \\
\midrule
Vomiting, diarrhea, heavy sweating & Rehydration requires appropriate water and electrolytes; concentrated products can be harmful.\\
Kidney, heart, or adrenal disease & Sodium, potassium, magnesium, and fluid handling may be altered.\\
Diuretics or blood-pressure medicines & Electrolytes can shift, sometimes without obvious symptoms.\\
\bottomrule
\end{tabular}
\end{table}

\section{Clinical assessment}
Ask about fluid intake, vomiting or diarrhea, sweating, eating disorders, alcohol, laxatives, diuretics, kidney and heart disease, and all prescribed and non-prescribed medicines. ECG assessment is important when potassium or magnesium disturbance is significant. In hospitalized patients, phosphate and magnesium may shift during nutritional rehabilitation; monitoring and gradual replacement are safer than indiscriminate dosing.
